\scanpage{110}
fibre, donc [les fibres étant intègres] $D'$ est transversal aux fibres. Enfin, \deleted{Et coïncide avec} par construction $D'_\eta=D_\eta$, donc $\Delta=D_{\deleted{\unsure{\eta}}}-D'_{\deleted{\unsure{\eta}}}$ est un diviseur tel que $\Delta_\eta=0$, $\operatorname{supp}\Delta$ ne contient aucune fibre. On voit facilement qu'un tel $\Delta$ est nul (\uline{lemme} : expliciter). On a donc $D=D'$, donc $D$ est Cartier.

\inserted{\uline{Cas général} Soit $S'$ l'ouvert de $X$ formé des pts où $f$ est simple. \deleted{Il suffit} Si $X\neq\varnothing$, alors $S'\longrightarrow S$ est simple, surjectif, donc on est ramené au cas de $X'/S'$.}
\margin{\unsure{Vu la normalité} de la base [expliciter]}

\uline{Question !} Supposons, au lieu de l'existence de $\underline{\operatorname{Pic}}_{X/S}$ et sa propreté essentielle, que le \uuline{foncteur} $\underline{\underline{\operatorname{Pic}}}_{X/S}$ ait la propriété de \deleted{prolongement} \inserted{propreté essentielle} [pour le \unsure{critère valuatif}], p.ex. $X\mathbin{-}S$ \uline{simple}, \deleted{ou $X=S\times_kY$, $Y/k$ normal}. \deleted{Est-il vrai alors} La conclusion du th. reste-t-elle valable ?
\margin{109}
